% GENERATED by build_do_now_tex.py on 2026-04-24 --- do not hand-edit; regenerate instead.
\documentclass[11pt]{article}
\usepackage{preamble}

\newcommand{\donowparabolagraph}{%
\begin{center}
\begin{tikzpicture}
\begin{axis}[
  width=0.82\linewidth,
  height=2.8in,
  xmin=-5.5,
  xmax=5.5,
  ymin=-18,
  ymax=28,
  axis lines=middle,
  xlabel={$x$},
  ylabel={$y$},
  xtick={-5,-4,-3,-2,-1,0,1,2,3,4,5},
  ytick={-16,-8,0,4,8,12,16,20,24},
  grid=major,
  tick label style={font=\small},
  label style={font=\small}
]
\addplot[tealheader, smooth, domain=-5:5, samples=250] {x^2};
\addplot[only marks, mark=*, mark size=2.2pt, color=dayblue] coordinates {
  (-4,16) (-3,9) (-2,4) (-1,1) (0,0) (1,1) (2,4) (3,9) (4,16) (5,25)
};
\node[anchor=south east, font=\small] at (axis cs:5,23) {$y=x^2$};
\end{axis}
\end{tikzpicture}
\end{center}
}

\begin{document}

\packetheading{Lesson 5-1 \textperiodcentered\ Period 1 \textperiodcentered\ Do Now}{L51 P1}

\vspace{0.25em}
\noindent\textbf{Name:}\ \rule{2in}{0.4pt}\quad
\textbf{Period:}\ \rule{0.5in}{0.4pt}\quad
\textbf{Date:}\ \rule{1in}{0.4pt}
\vspace{0.5em}

\bankitem{Do Now --- Explore \& Reason: \(y = x^2\) graph}{%
%
\textbf{EXPLORE \& REASON}

The graph shows \(y = x^2\).

\textit{[GRAPH / TIKZ figure]}

\donowparabolagraph

\textbf{A.} Find all possible values of \(x\) or \(y\) so that the point is on the graph.
\begin{enumerate}[label=(\alph*),leftmargin=2.2em,itemsep=2pt,topsep=3pt]
  \item \((2, \underline{\hspace{1.2cm}})\)
  \item \((3, \underline{\hspace{1.2cm}})\)
  \item \((-3, \underline{\hspace{1.2cm}})\)
  \item \((5, \underline{\hspace{1.2cm}})\)
  \item \((\underline{\hspace{1.2cm}}, 4)\)
  \item \((\underline{\hspace{1.2cm}}, -16)\)
  \item \((\underline{\hspace{1.2cm}}, 7)\)
  \item \((\underline{\hspace{1.2cm}}, 5)\)
\end{enumerate}

\textbf{B. Communicate Precisely} Write a precise set of instructions that show how to find an approximate value of \(\sqrt{13}\) using the graph.

\textbf{C.} Draw a graph of \(y = x^3\). Use the graph to approximate each value.
\begin{enumerate}[label=(\alph*),leftmargin=2.2em,itemsep=2pt,topsep=3pt]
  \item \(\sqrt[3]{5}\)
  \item \(\sqrt[3]{-5}\)
  \item \(\sqrt[3]{8}\)
  \item A solution to \(x^3 = 5\)
  \item A solution to \(x^3 = -5\)
  \item A solution to \(x^3 = 8\)
\end{enumerate}
}{}
\vspace{2.5in}

\end{document}
